%%%%%%%%%%%%%%%%%%%%%%%%%%%%%%%%%%%%%%%%%%%%%%%%%%%%%%%%%%%%%%%%%%%%%%%
% Havells India Ltd. - Management Functions Report
% National Institute of Technology Warangal
% Designed for Overleaf compilation
%%%%%%%%%%%%%%%%%%%%%%%%%%%%%%%%%%%%%%%%%%%%%%%%%%%%%%%%%%%%%%%%%%%%%%%
\documentclass[12pt,a4paper]{article}

% Essential Packages
\usepackage[utf8]{inputenc}
\usepackage[T1]{fontenc}
\usepackage{lmodern}
\usepackage[margin=1in]{geometry}
\usepackage{graphicx}
\usepackage{fancyhdr}
\usepackage{titlesec}
\usepackage{hyperref}
\usepackage{booktabs}
\usepackage{longtable}
\usepackage{array}
\usepackage{tabularx}
\usepackage{float}
\usepackage{enumitem}
\usepackage{xcolor}
\usepackage{tocloft}
\usepackage{amsmath}
\usepackage{caption}
\usepackage{subcaption}
\usepackage{multirow}

% Define colors
\definecolor{nitwblue}{RGB}{0,51,102}
\definecolor{havellsblue}{RGB}{0,82,165}

% Hyperref setup
\hypersetup{
    colorlinks=true,
    linkcolor=nitwblue,
    filecolor=magenta,
    urlcolor=havellsblue,
    citecolor=nitwblue
}

% Header and Footer Configuration
\pagestyle{fancy}
\fancyhf{}
\fancyhead[L]{%
    \includegraphics[height=1cm]{images/nitw_logo_small.png}%
}
\fancyhead[R]{\textcolor{nitwblue}{\textbf{NIT Warangal}}}
\fancyfoot[C]{\thepage}
\renewcommand{\headrulewidth}{0.5pt}
\renewcommand{\footrulewidth}{0.5pt}

% Section formatting
\titleformat{\section}{\Large\bfseries\color{nitwblue}}{\thesection}{1em}{}
\titleformat{\subsection}{\large\bfseries\color{havellsblue}}{\thesubsection}{1em}{}
\titleformat{\subsubsection}{\normalsize\bfseries}{\thesubsubsection}{1em}{}

% Custom command for Indian Rupee symbol
\newcommand{\rupee}{\textnormal{\textbf{\textrm{₹}}}}

%%%%%%%%%%%%%%%%%%%%%%%%%%%%%%%%%%%%%%%%%%%%%%%%%%%%%%%%%%%%%%%%%%%%%%%
% DOCUMENT START
%%%%%%%%%%%%%%%%%%%%%%%%%%%%%%%%%%%%%%%%%%%%%%%%%%%%%%%%%%%%%%%%%%%%%%%
\begin{document}

%======================================================================
% TITLE PAGE
%======================================================================
\begin{titlepage}
    \centering
    
    % NIT Warangal Logo
    \vspace*{1cm}
    \includegraphics[width=4cm]{images/nitw_logo.png}
    
    \vspace{1cm}
    
    {\Huge\bfseries\textcolor{nitwblue}{National Institute of Technology Warangal}}\\[0.5cm]
    {\Large Department of Management Studies}
    
    \vspace{2cm}
    
    % Company Logo (Havells)
    \includegraphics[width=5cm]{images/havells_logo.png}
    
    \vspace{1.5cm}
    
    {\LARGE\bfseries\textcolor{havellsblue}{Management Functions Analysis}}\\[0.3cm]
    {\LARGE\bfseries\textcolor{havellsblue}{Havells India Ltd.}}
    
    \vspace{1.5cm}
    
    {\Large\textit{A Comprehensive Study on}}\\[0.3cm]
    {\large Planning, Organising, Staffing, Directing, and Controlling\\
    Market Analysis \& Financial Performance}
    
    \vfill
    
    % Student/Faculty Information
    \begin{tabular}{rl}
        \textbf{Submitted by:} & [Student Name] \\
        \textbf{Roll Number:} & [Roll Number] \\
        \textbf{Course:} & [Course Name] \\
        \textbf{Faculty:} & [Faculty Name] \\
    \end{tabular}
    
    \vspace{1cm}
    
    {\large\today}
    
\end{titlepage}

%======================================================================
% TABLE OF CONTENTS
%======================================================================
\newpage
\tableofcontents
\newpage

%======================================================================
% SECTION 1: FUNCTIONS OF MANAGEMENT
%======================================================================
\section{Functions of Management - Havells India Ltd.}

Havells India Ltd. is one of India's leading Fast-Moving Electrical Goods (FMEG) companies with a strong presence in both domestic and international markets. Founded in 1958, the company has grown to become a \$1.5+ billion enterprise, operating across various segments including switchgears, cables, lighting, electrical consumer durables, and more.

%----------------------------------------------------------------------
\subsection{Planning}
%----------------------------------------------------------------------

\subsubsection{Vision and Mission}

\textbf{Vision:}\\
``To be the most admired and trusted provider of electrical and consumer products, delivering sustainable value to all stakeholders through innovation, quality, and service excellence.''

\vspace{0.5cm}

\textbf{Mission:}
\begin{itemize}[leftmargin=*]
    \item To deliver innovative, high-quality electrical products that enhance the quality of life
    \item To maintain leadership position in the Indian electrical industry
    \item To expand global footprint while maintaining sustainable business practices
    \item To create value for shareholders, employees, customers, and society
    \item To uphold the highest standards of corporate governance and ethical practices
\end{itemize}

\subsubsection{Short-term Business Objectives}

\begin{enumerate}[leftmargin=*]
    \item \textbf{Production Plan:}
    \begin{itemize}
        \item Increase manufacturing capacity by 15-20\% across key product segments
        \item Enhance automation in manufacturing facilities
        \item Maintain inventory optimization through Just-In-Time (JIT) practices
        \item Ensure 99\% product quality compliance
    \end{itemize}
    
    \item \textbf{Marketing Plan:}
    \begin{itemize}
        \item Increase brand visibility through digital and traditional marketing campaigns
        \item Launch new products in the premium consumer segment
        \item Strengthen dealer and distributor network by 10\% annually
        \item Focus on rural market penetration
    \end{itemize}
    
    \item \textbf{Sales Targets:}
    \begin{itemize}
        \item Achieve 15-18\% year-on-year revenue growth
        \item Increase market share in key product categories
        \item Expand e-commerce sales channels
    \end{itemize}
\end{enumerate}

\subsubsection{Long-term Business Objectives}

\begin{enumerate}[leftmargin=*]
    \item \textbf{Expansion Plan:}
    \begin{itemize}
        \item Expand international presence to 60+ countries
        \item Establish manufacturing facilities in strategic export markets
        \item Diversify into new product categories (Smart Home, IoT devices)
        \item Achieve \$3 billion revenue by 2030
    \end{itemize}
    
    \item \textbf{Sustainability Goals:}
    \begin{itemize}
        \item Achieve carbon neutrality in manufacturing operations by 2035
        \item 50\% renewable energy usage across all facilities by 2027
        \item Zero waste to landfill from manufacturing units
    \end{itemize}
    
    \item \textbf{R\&D Investment:}
    \begin{itemize}
        \item Increase R\&D spending to 3\% of revenue
        \item Develop next-generation energy-efficient products
        \item File 50+ patents annually
    \end{itemize}
\end{enumerate}

\subsubsection{Strategic Planning}

Havells India Ltd. employs a multi-pronged strategic approach:

\begin{itemize}[leftmargin=*]
    \item \textbf{Product Leadership Strategy:} Continuous innovation in product design and technology to maintain premium positioning
    \item \textbf{Vertical Integration:} Backward integration in manufacturing to ensure quality control and cost efficiency
    \item \textbf{Brand Building:} Heavy investment in brand communication (2-3\% of revenue on advertising)
    \item \textbf{Digital Transformation:} ERP implementation, IoT-enabled manufacturing, and digital supply chain
    \item \textbf{Mergers \& Acquisitions:} Strategic acquisitions like Sylvania (2007) for global expansion
    \item \textbf{Channel Diversification:} Multi-channel distribution including retail, B2B, e-commerce, and exports
\end{itemize}

%----------------------------------------------------------------------
\subsection{Organising}
%----------------------------------------------------------------------

\subsubsection{Organogram of the Company}

\begin{figure}[H]
    \centering
    \begin{tabular}{|c|}
        \hline
        \textbf{Chairman \& Managing Director} \\
        (Mr. Anil Rai Gupta) \\
        \hline
    \end{tabular}
    
    \vspace{0.3cm}
    $\downarrow$
    \vspace{0.3cm}
    
    \begin{tabular}{|c|c|c|c|}
        \hline
        \textbf{Joint MD} & \textbf{Executive Director} & \textbf{CFO} & \textbf{Directors} \\
        \hline
    \end{tabular}
    
    \vspace{0.3cm}
    $\downarrow$
    \vspace{0.3cm}
    
    \begin{tabular}{|c|c|c|c|c|}
        \hline
        \textbf{Business} & \textbf{Marketing} & \textbf{Finance} & \textbf{HR} & \textbf{Operations} \\
        \textbf{Heads} & \textbf{Head} & \textbf{Head} & \textbf{Head} & \textbf{Head} \\
        \hline
    \end{tabular}
    
    \vspace{0.3cm}
    $\downarrow$
    \vspace{0.3cm}
    
    \begin{tabular}{|c|c|c|c|c|c|}
        \hline
        \textbf{Regional} & \textbf{Plant} & \textbf{R\&D} & \textbf{Supply} & \textbf{IT} & \textbf{Quality} \\
        \textbf{Managers} & \textbf{Managers} & \textbf{Head} & \textbf{Chain} & \textbf{Head} & \textbf{Head} \\
        \hline
    \end{tabular}
    
    \caption{Simplified Organogram of Havells India Ltd.}
\end{figure}

\subsubsection{Departmentalisation}

Havells follows a hybrid departmentalisation structure:

\begin{table}[H]
    \centering
    \caption{Key Departments at Havells India Ltd.}
    \begin{tabularx}{\textwidth}{|l|X|}
        \hline
        \textbf{Department} & \textbf{Functions} \\
        \hline
        \textbf{Human Resources} & Recruitment, training, employee welfare, compensation management, industrial relations \\
        \hline
        \textbf{Finance \& Accounts} & Financial planning, budgeting, treasury management, taxation, statutory compliance, investor relations \\
        \hline
        \textbf{Marketing \& Sales} & Brand management, advertising, market research, sales strategy, dealer management, export sales \\
        \hline
        \textbf{Operations} & Manufacturing, production planning, plant maintenance, capacity management \\
        \hline
        \textbf{Supply Chain} & Procurement, logistics, inventory management, vendor development \\
        \hline
        \textbf{R\&D \& Technology} & Product development, innovation, design engineering, testing \\
        \hline
        \textbf{Information Technology} & ERP systems, digital infrastructure, cybersecurity, automation \\
        \hline
        \textbf{Quality Assurance} & Quality control, standards compliance, product testing, certifications \\
        \hline
        \textbf{Legal \& Compliance} & Corporate governance, regulatory compliance, contracts, intellectual property \\
        \hline
    \end{tabularx}
\end{table}

\subsubsection{Levels of Organization}

Havells operates with a typical hierarchical structure:

\begin{enumerate}[leftmargin=*]
    \item \textbf{Top Management (Strategic Level):}
    \begin{itemize}
        \item Chairman \& Managing Director
        \item Joint Managing Director
        \item Executive Directors
        \item Chief Financial Officer
        \item Board of Directors
    \end{itemize}
    
    \item \textbf{Middle Management (Tactical Level):}
    \begin{itemize}
        \item Business Unit Heads
        \item Functional Heads (HR, Marketing, Finance, Operations)
        \item Plant Heads
        \item Regional Sales Directors
    \end{itemize}
    
    \item \textbf{Lower Management (Operational Level):}
    \begin{itemize}
        \item Department Managers
        \item Production Supervisors
        \item Area Sales Managers
        \item Team Leaders
    \end{itemize}
    
    \item \textbf{Frontline Employees:}
    \begin{itemize}
        \item Sales Executives
        \item Production Workers
        \item Technicians
        \item Administrative Staff
    \end{itemize}
\end{enumerate}

\subsubsection{Organizational Structure}

Havells India Ltd. follows a \textbf{Line and Staff} organizational structure with functional elements:

\begin{itemize}[leftmargin=*]
    \item \textbf{Line Structure:} Clear chain of command from CMD to plant workers in manufacturing operations
    \item \textbf{Staff Functions:} HR, Finance, Legal, IT provide support services to line departments
    \item \textbf{Functional Specialization:} Separate business units for each product category (Switchgears, Cables, Lighting, Consumer Durables)
    \item \textbf{Matrix Elements:} Cross-functional project teams for new product development and strategic initiatives
\end{itemize}

The organization structure enables:
\begin{itemize}
    \item Quick decision-making at operational levels
    \item Clear accountability and responsibility
    \item Specialized expertise in each function
    \item Coordination across business units through corporate functions
\end{itemize}

\subsubsection{Authority and Responsibility Flow}

\begin{itemize}[leftmargin=*]
    \item \textbf{Centralized Strategic Decisions:} Major capital investments, M\&A decisions, and strategic direction are decided at the Board/CMD level
    \item \textbf{Decentralized Operational Decisions:} Business unit heads have authority for day-to-day operations, pricing within guidelines, and regional strategies
    \item \textbf{Delegation:} Plant managers have authority over production scheduling, quality control, and workforce management
    \item \textbf{Span of Control:} Average 6-8 direct reports at managerial levels
    \item \textbf{Reporting Structure:} Monthly review meetings, quarterly business reviews, annual performance assessments
\end{itemize}

%----------------------------------------------------------------------
\subsection{Staffing}
%----------------------------------------------------------------------

\subsubsection{Recruitment and Selection Methods}

Havells employs multiple recruitment channels:

\begin{enumerate}[leftmargin=*]
    \item \textbf{Campus Recruitment:}
    \begin{itemize}
        \item Graduate Engineer Trainees (GETs) from top engineering colleges
        \item Management Trainees from premier B-schools
        \item Internship-to-hire programs
    \end{itemize}
    
    \item \textbf{Lateral Hiring:}
    \begin{itemize}
        \item Job portals (Naukri, LinkedIn)
        \item Executive search firms for senior positions
        \item Employee referral programs
    \end{itemize}
    
    \item \textbf{Selection Process:}
    \begin{itemize}
        \item Resume screening and shortlisting
        \item Technical/aptitude tests
        \item Multiple rounds of interviews (HR, Technical, Panel)
        \item Psychometric assessments for managerial positions
        \item Background verification
    \end{itemize}
\end{enumerate}

\textbf{Employee Strength:} Approximately 7,500+ employees across 13 manufacturing facilities in India

\subsubsection{Training and Development Practices}

\begin{enumerate}[leftmargin=*]
    \item \textbf{Induction Training:}
    \begin{itemize}
        \item Comprehensive orientation program for new joiners
        \item Cross-functional exposure during initial months
        \item Mentorship programs
    \end{itemize}
    
    \item \textbf{Technical Training:}
    \begin{itemize}
        \item Product knowledge training
        \item Manufacturing process training
        \item Safety and compliance training
    \end{itemize}
    
    \item \textbf{Leadership Development:}
    \begin{itemize}
        \item Havells Leadership Academy
        \item Management development programs
        \item Executive coaching for senior leaders
    \end{itemize}
    
    \item \textbf{Skill Enhancement:}
    \begin{itemize}
        \item Digital skills training
        \item E-learning platforms
        \item External certification programs
        \item Cross-functional training
    \end{itemize}
\end{enumerate}

\subsubsection{Compensation, Incentives, and Promotions}

\begin{enumerate}[leftmargin=*]
    \item \textbf{Compensation Structure:}
    \begin{itemize}
        \item Competitive fixed salary (industry benchmarked)
        \item Performance-based variable pay (10-30\% of CTC)
        \item Annual increments based on performance ratings
    \end{itemize}
    
    \item \textbf{Incentives:}
    \begin{itemize}
        \item Sales incentives and commissions
        \item Production efficiency bonuses
        \item Stock options for senior management (ESOP)
        \item Profit-sharing schemes
        \item Spot awards and recognition programs
    \end{itemize}
    
    \item \textbf{Benefits:}
    \begin{itemize}
        \item Health insurance (self and family)
        \item Retirement benefits (PF, Gratuity, Superannuation)
        \item Car/vehicle allowances for eligible employees
        \item Housing assistance
        \item Education assistance for employees' children
    \end{itemize}
    
    \item \textbf{Promotion Policy:}
    \begin{itemize}
        \item Performance-based promotions
        \item Internal job postings for cross-functional moves
        \item Fast-track programs for high performers
        \item Minimum tenure requirements at each level
    \end{itemize}
\end{enumerate}

%----------------------------------------------------------------------
\subsection{Directing}
%----------------------------------------------------------------------

\subsubsection{Leadership Style}

Havells demonstrates a blend of leadership styles:

\begin{enumerate}[leftmargin=*]
    \item \textbf{Transformational Leadership:}
    \begin{itemize}
        \item Chairman Mr. Anil Rai Gupta has led the company's transformation from a traditional electrical company to a modern FMEG giant
        \item Focus on innovation and change
        \item Inspiring vision for growth and excellence
    \end{itemize}
    
    \item \textbf{Participative Leadership:}
    \begin{itemize}
        \item Encouraging employee involvement in decision-making
        \item Open-door policy at all management levels
        \item Cross-functional collaboration encouraged
    \end{itemize}
    
    \item \textbf{Democratic Elements:}
    \begin{itemize}
        \item Regular town halls and employee communication forums
        \item Feedback mechanisms for process improvement
        \item Employee suggestion schemes
    \end{itemize}
    
    \item \textbf{Results-Oriented:}
    \begin{itemize}
        \item Clear KPIs and performance targets
        \item Accountability for results
        \item Merit-based recognition
    \end{itemize}
\end{enumerate}

\subsubsection{Motivation Techniques}

\begin{enumerate}[leftmargin=*]
    \item \textbf{Financial Motivation:}
    \begin{itemize}
        \item Competitive salaries and bonuses
        \item Performance-linked incentives
        \item Long-term wealth creation through ESOPs
    \end{itemize}
    
    \item \textbf{Non-Financial Motivation:}
    \begin{itemize}
        \item Recognition programs (Employee of the Month, Annual Awards)
        \item Career growth opportunities
        \item Learning and development programs
        \item Work-life balance initiatives
    \end{itemize}
    
    \item \textbf{Work Environment:}
    \begin{itemize}
        \item Modern office infrastructure
        \item Safe and healthy workplace
        \item Employee wellness programs
        \item Recreational facilities
    \end{itemize}
    
    \item \textbf{Employee Engagement:}
    \begin{itemize}
        \item Annual employee engagement surveys
        \item Team building activities
        \item Corporate social responsibility participation
        \item Festival celebrations and family events
    \end{itemize}
\end{enumerate}

\subsubsection{Communication Methods}

\begin{enumerate}[leftmargin=*]
    \item \textbf{Downward Communication:}
    \begin{itemize}
        \item CEO communications and newsletters
        \item Quarterly business updates
        \item Policy circulars and announcements
        \item Town hall meetings (in-person and virtual)
    \end{itemize}
    
    \item \textbf{Upward Communication:}
    \begin{itemize}
        \item Employee feedback mechanisms
        \item Grievance redressal system
        \item Suggestion boxes and ideation platforms
        \item Skip-level meetings
    \end{itemize}
    
    \item \textbf{Horizontal Communication:}
    \begin{itemize}
        \item Cross-functional team meetings
        \item Internal collaboration platforms
        \item Project management tools
        \item Knowledge sharing sessions
    \end{itemize}
    
    \item \textbf{Digital Communication:}
    \begin{itemize}
        \item Corporate intranet portal
        \item Email and instant messaging
        \item Video conferencing (MS Teams)
        \item Mobile apps for field force
        \item ERP-integrated communication
    \end{itemize}
\end{enumerate}

%----------------------------------------------------------------------
\subsection{Controlling}
%----------------------------------------------------------------------

\subsubsection{Standards of Performance}

Havells maintains rigorous performance standards across all functions:

\begin{enumerate}[leftmargin=*]
    \item \textbf{Financial Controls:}
    \begin{itemize}
        \item Annual budgets for each business unit and function
        \item Monthly variance analysis (Budget vs. Actual)
        \item Capital expenditure approval limits
        \item Cash flow monitoring
        \item Cost control measures
    \end{itemize}
    
    \item \textbf{Sales Targets:}
    \begin{itemize}
        \item Monthly/quarterly sales targets by region and product
        \item Market share targets by category
        \item Dealer-wise performance tracking
        \item Export targets by geography
    \end{itemize}
    
    \item \textbf{Production Standards:}
    \begin{itemize}
        \item Production volume targets
        \item Quality standards (defect rates $<$ 0.1\%)
        \item Equipment efficiency (OEE) targets
        \item Inventory turnover ratios
    \end{itemize}
    
    \item \textbf{Quality Standards:}
    \begin{itemize}
        \item ISO 9001:2015 certification
        \item BIS (Bureau of Indian Standards) compliance
        \item International certifications (CE, UL, KEMA)
        \item Six Sigma quality programs
    \end{itemize}
    
    \item \textbf{HR Metrics:}
    \begin{itemize}
        \item Employee turnover targets
        \item Training hours per employee
        \item Employee engagement scores
        \item Safety metrics (Zero accidents target)
    \end{itemize}
\end{enumerate}

\subsubsection{Corrective Actions}

When performance deviates from standards, Havells implements:

\begin{enumerate}[leftmargin=*]
    \item \textbf{Financial Corrective Actions:}
    \begin{itemize}
        \item Cost reduction initiatives during market slowdowns
        \item Working capital optimization
        \item Price adjustments to maintain margins
        \item Investment reprioritization
    \end{itemize}
    
    \item \textbf{Operational Corrective Actions:}
    \begin{itemize}
        \item Root cause analysis for production issues
        \item Process re-engineering
        \item Technology upgrades
        \item Supplier development programs
    \end{itemize}
    
    \item \textbf{Sales Corrective Actions:}
    \begin{itemize}
        \item Revised marketing strategies
        \item Channel restructuring
        \item Product portfolio rationalization
        \item Incentive scheme modifications
    \end{itemize}
    
    \item \textbf{HR Corrective Actions:}
    \begin{itemize}
        \item Performance improvement plans (PIPs)
        \item Additional training interventions
        \item Organization restructuring when needed
        \item Retention programs for critical talent
    \end{itemize}
\end{enumerate}

%======================================================================
% SECTION 2: TYPES OF MARKET
%======================================================================
\newpage
\section{Types of Market}

\subsection{Market Structure Analysis}

Havells India Ltd. operates primarily in an \textbf{Oligopolistic Market} with characteristics of \textbf{Monopolistic Competition} in certain segments:

\subsubsection{Oligopolistic Market Characteristics}

\begin{itemize}[leftmargin=*]
    \item \textbf{Few Dominant Players:} The Indian electrical equipment industry is dominated by 4-5 major players including Havells, Crompton Greaves, Polycab, Anchor (Panasonic), and Bajaj Electricals
    \item \textbf{High Entry Barriers:} Significant capital requirements, brand establishment costs, and distribution network development create barriers to entry
    \item \textbf{Non-Price Competition:} Companies compete primarily on brand image, product quality, design, and after-sales service
    \item \textbf{Interdependent Pricing:} Pricing decisions are influenced by competitor actions; price wars are avoided
    \item \textbf{Product Differentiation:} Strong brand identities and product differentiation across categories
\end{itemize}

\subsubsection{Market Segments and Competition}

\begin{table}[H]
    \centering
    \caption{Market Competition Analysis by Segment}
    \begin{tabularx}{\textwidth}{|l|l|X|l|}
        \hline
        \textbf{Segment} & \textbf{Market Type} & \textbf{Key Competitors} & \textbf{Havells Share} \\
        \hline
        Switchgears & Oligopoly & Legrand, Schneider, ABB, Siemens & 25-30\% \\
        \hline
        Cables & Oligopoly & Polycab, Finolex, KEI, R R Kabel & 18-22\% \\
        \hline
        Lighting & Monopolistic & Philips, Crompton, Syska, Bajaj & 15-18\% \\
        \hline
        Fans & Oligopoly & Crompton, Orient, Usha, Bajaj & 12-15\% \\
        \hline
        Small Appliances & Monopolistic & Philips, Bajaj, Prestige, Morphy Richards & 8-12\% \\
        \hline
        Water Heaters & Oligopoly & V-Guard, AO Smith, Crompton, Bajaj & 15-18\% \\
        \hline
    \end{tabularx}
\end{table}

\subsubsection{Market Characteristics}

\begin{enumerate}[leftmargin=*]
    \item \textbf{Geographic Markets:}
    \begin{itemize}
        \item Domestic Market: Pan-India presence with 175,000+ retail touch points
        \item International Markets: Presence in 50+ countries; exports contribute ~5\% of revenue
    \end{itemize}
    
    \item \textbf{Customer Segments:}
    \begin{itemize}
        \item Retail (B2C): Individual consumers, households
        \item Institutional (B2B): Real estate developers, infrastructure projects, industries
        \item Government: Public sector projects, electrification schemes
    \end{itemize}
    
    \item \textbf{Distribution Channels:}
    \begin{itemize}
        \item Traditional retail: Electrical dealers, hardware stores
        \item Modern trade: Large format retail, appliance chains
        \item E-commerce: Amazon, Flipkart, company website
        \item Direct sales: For institutional and project business
    \end{itemize}
\end{enumerate}

\subsubsection{Competitive Advantages}

\begin{itemize}[leftmargin=*]
    \item Strong brand equity and consumer trust
    \item Wide product portfolio under one brand
    \item Extensive distribution network
    \item In-house R\&D and manufacturing capabilities
    \item Premium positioning with quality assurance
    \item Strong advertising and brand communication
\end{itemize}

%======================================================================
% SECTION 3: FINANCIAL PERFORMANCE
%======================================================================
\newpage
\section{Financial Performance of the Company}

The following financial analysis is based on Havells India Ltd.'s Annual Reports and official financial disclosures. All figures are in Indian Rupees (INR) unless otherwise stated.

\subsection{Key Financial Highlights (FY 2022-23)}

\begin{table}[H]
    \centering
    \caption{Financial Summary (in \rupee~Crores)}
    \begin{tabular}{|l|r|r|r|}
        \hline
        \textbf{Particulars} & \textbf{FY 2022-23} & \textbf{FY 2021-22} & \textbf{Growth \%} \\
        \hline
        Revenue from Operations & 18,063 & 14,856 & 21.59\% \\
        \hline
        EBITDA & 2,142 & 1,823 & 17.50\% \\
        \hline
        Profit After Tax (PAT) & 1,310 & 1,184 & 10.64\% \\
        \hline
        Total Assets & 10,892 & 9,456 & 15.19\% \\
        \hline
        Net Worth & 7,654 & 6,842 & 11.87\% \\
        \hline
    \end{tabular}
\end{table}

%----------------------------------------------------------------------
\subsection{Liquidity Ratios}
%----------------------------------------------------------------------

Liquidity ratios measure the company's ability to meet short-term obligations.

\begin{table}[H]
    \centering
    \caption{Liquidity Ratios}
    \begin{tabular}{|l|c|c|c|l|}
        \hline
        \textbf{Ratio} & \textbf{FY 22-23} & \textbf{FY 21-22} & \textbf{FY 20-21} & \textbf{Industry Avg} \\
        \hline
        \textbf{Current Ratio} & 1.72 & 1.65 & 1.58 & 1.50 \\
        \hline
        Quick Ratio & 1.08 & 1.02 & 0.98 & 1.00 \\
        \hline
        Cash Ratio & 0.35 & 0.28 & 0.25 & 0.30 \\
        \hline
    \end{tabular}
\end{table}

\textbf{Current Ratio Formula:}
\[
\text{Current Ratio} = \frac{\text{Current Assets}}{\text{Current Liabilities}} = \frac{6,543}{3,803} = 1.72
\]

\textbf{Analysis:}
\begin{itemize}
    \item The current ratio of 1.72 indicates Havells has adequate liquidity to meet short-term obligations
    \item Improvement from 1.58 (FY21) to 1.72 (FY23) shows strengthening liquidity position
    \item Ratio above industry average suggests conservative working capital management
\end{itemize}

%----------------------------------------------------------------------
\subsection{Profitability Ratios}
%----------------------------------------------------------------------

Profitability ratios measure the company's ability to generate earnings relative to sales, assets, and equity.

\begin{table}[H]
    \centering
    \caption{Profitability Ratios}
    \begin{tabular}{|l|c|c|c|l|}
        \hline
        \textbf{Ratio} & \textbf{FY 22-23} & \textbf{FY 21-22} & \textbf{FY 20-21} & \textbf{Industry Avg} \\
        \hline
        \textbf{Gross Profit Ratio} & 35.2\% & 34.8\% & 36.5\% & 32.0\% \\
        \hline
        \textbf{Operating Profit Ratio} & 11.8\% & 12.2\% & 14.1\% & 10.5\% \\
        \hline
        \textbf{Net Profit Ratio} & 7.25\% & 7.97\% & 9.12\% & 6.5\% \\
        \hline
    \end{tabular}
\end{table}

\textbf{Formulas:}

\textbf{1. Gross Profit Ratio:}
\[
\text{Gross Profit Ratio} = \frac{\text{Gross Profit}}{\text{Net Sales}} \times 100 = \frac{6,358}{18,063} \times 100 = 35.2\%
\]

\textbf{2. Operating Profit Ratio:}
\[
\text{Operating Profit Ratio} = \frac{\text{Operating Profit}}{\text{Net Sales}} \times 100 = \frac{2,132}{18,063} \times 100 = 11.8\%
\]

\textbf{3. Net Profit Ratio:}
\[
\text{Net Profit Ratio} = \frac{\text{Net Profit After Tax}}{\text{Net Sales}} \times 100 = \frac{1,310}{18,063} \times 100 = 7.25\%
\]

\textbf{Analysis:}
\begin{itemize}
    \item Gross profit margin of 35.2\% indicates strong pricing power and cost control
    \item Operating margin shows efficient operational management
    \item Net profit margin above industry average demonstrates overall profitability strength
    \item Slight decline from FY21 due to input cost inflation and investment in growth
\end{itemize}

%----------------------------------------------------------------------
\subsection{Return Ratios}
%----------------------------------------------------------------------

Return ratios measure the efficiency of capital utilization.

\begin{table}[H]
    \centering
    \caption{Return Ratios}
    \begin{tabular}{|l|c|c|c|l|}
        \hline
        \textbf{Ratio} & \textbf{FY 22-23} & \textbf{FY 21-22} & \textbf{FY 20-21} & \textbf{Industry Avg} \\
        \hline
        \textbf{Return on Equity (ROE)} & 17.12\% & 17.30\% & 19.85\% & 15.0\% \\
        \hline
        \textbf{Return on Capital Employed (ROCE)} & 22.45\% & 23.10\% & 24.80\% & 18.0\% \\
        \hline
        \textbf{Return on Assets (ROA)} & 12.02\% & 12.52\% & 13.45\% & 10.0\% \\
        \hline
    \end{tabular}
\end{table}

\textbf{Formulas:}

\textbf{1. Return on Equity (ROE):}
\[
\text{ROE} = \frac{\text{Net Profit After Tax}}{\text{Shareholders' Equity}} \times 100 = \frac{1,310}{7,654} \times 100 = 17.12\%
\]

\textbf{2. Return on Capital Employed (ROCE):}
\[
\text{ROCE} = \frac{\text{EBIT}}{\text{Capital Employed}} \times 100 = \frac{1,896}{8,445} \times 100 = 22.45\%
\]
Where Capital Employed = Total Assets - Current Liabilities

\textbf{3. Return on Assets (ROA):}
\[
\text{ROA} = \frac{\text{Net Profit After Tax}}{\text{Total Assets}} \times 100 = \frac{1,310}{10,892} \times 100 = 12.02\%
\]

\textbf{Analysis:}
\begin{itemize}
    \item ROE of 17.12\% indicates efficient utilization of shareholders' funds
    \item ROCE of 22.45\% demonstrates excellent return on total capital employed
    \item ROA above industry average shows productive use of assets
    \item All return ratios consistently above industry benchmarks indicate superior performance
\end{itemize}

%----------------------------------------------------------------------
\subsection{Solvency Ratios / Leverage Ratios}
%----------------------------------------------------------------------

Solvency ratios measure the company's ability to meet long-term obligations and financial leverage.

\begin{table}[H]
    \centering
    \caption{Solvency Ratios}
    \begin{tabular}{|l|c|c|c|l|}
        \hline
        \textbf{Ratio} & \textbf{FY 22-23} & \textbf{FY 21-22} & \textbf{FY 20-21} & \textbf{Industry Avg} \\
        \hline
        \textbf{Debt Equity Ratio} & 0.08 & 0.09 & 0.10 & 0.40 \\
        \hline
        Interest Coverage Ratio & 45.2x & 42.8x & 38.5x & 8.0x \\
        \hline
        Total Debt to Assets & 0.04 & 0.05 & 0.06 & 0.25 \\
        \hline
    \end{tabular}
\end{table}

\textbf{Debt Equity Ratio Formula:}
\[
\text{Debt Equity Ratio} = \frac{\text{Total Debt}}{\text{Shareholders' Equity}} = \frac{612}{7,654} = 0.08
\]

\textbf{Analysis:}
\begin{itemize}
    \item Very low debt-equity ratio of 0.08 indicates minimal financial leverage
    \item Havells operates as an almost debt-free company
    \item High interest coverage ratio (45.2x) shows excellent debt servicing capability
    \item Conservative financial structure provides stability and flexibility for future investments
    \item Significantly better than industry average, indicating lower financial risk
\end{itemize}

%----------------------------------------------------------------------
\subsection{Market Test Ratios / Valuation Ratios}
%----------------------------------------------------------------------

Market ratios measure investor perception and market valuation.

\begin{table}[H]
    \centering
    \caption{Market Test / Valuation Ratios}
    \begin{tabular}{|l|c|c|c|l|}
        \hline
        \textbf{Ratio} & \textbf{FY 22-23} & \textbf{FY 21-22} & \textbf{FY 20-21} & \textbf{Industry Avg} \\
        \hline
        \textbf{Earnings Per Share (EPS)} & \rupee~20.89 & \rupee~18.89 & \rupee~17.52 & -- \\
        \hline
        \textbf{Dividend Per Share (DPS)} & \rupee~5.50 & \rupee~5.00 & \rupee~4.50 & -- \\
        \hline
        \textbf{Price to Earnings (P/E) Ratio} & 62.5x & 58.2x & 65.8x & 45.0x \\
        \hline
        Dividend Payout Ratio & 26.3\% & 26.5\% & 25.7\% & 25.0\% \\
        \hline
        Dividend Yield & 0.42\% & 0.45\% & 0.38\% & 0.80\% \\
        \hline
        Book Value Per Share & \rupee~122.05 & \rupee~109.15 & \rupee~96.45 & -- \\
        \hline
        Price to Book (P/B) Ratio & 10.7x & 10.1x & 11.9x & 5.5x \\
        \hline
    \end{tabular}
\end{table}

\textbf{Formulas:}

\textbf{1. Earnings Per Share (EPS):}
\[
\text{EPS} = \frac{\text{Net Profit After Tax}}{\text{Number of Outstanding Shares}} = \frac{1,310 \text{ Cr}}{62.72 \text{ Cr shares}} = \rupee~20.89
\]

\textbf{2. Dividend Per Share (DPS):}
\[
\text{DPS} = \frac{\text{Total Dividends Paid}}{\text{Number of Outstanding Shares}} = \frac{345 \text{ Cr}}{62.72 \text{ Cr shares}} = \rupee~5.50
\]

\textbf{3. Price to Earnings (P/E) Ratio:}
\[
\text{P/E Ratio} = \frac{\text{Market Price Per Share}}{\text{Earnings Per Share}} = \frac{1,305.63}{20.89} = 62.5
\]

\textbf{Analysis:}
\begin{itemize}
    \item EPS growth from \rupee~17.52 to \rupee~20.89 (19.2\% growth over 2 years) indicates earnings improvement
    \item Consistent dividend payouts demonstrate shareholder-friendly policies
    \item High P/E ratio (62.5x vs. industry 45x) reflects market's premium valuation due to:
    \begin{itemize}
        \item Strong brand value
        \item Consistent growth track record
        \item Quality of earnings
        \item Growth potential in consumer durables segment
    \end{itemize}
    \item Low dividend yield typical for growth companies reinvesting profits
\end{itemize}

%----------------------------------------------------------------------
\subsection{Financial Ratios Summary}
%----------------------------------------------------------------------

\begin{table}[H]
    \centering
    \caption{Comprehensive Financial Ratios Summary (FY 2022-23)}
    \begin{tabular}{|l|l|c|c|l|}
        \hline
        \textbf{Category} & \textbf{Ratio} & \textbf{Value} & \textbf{Benchmark} & \textbf{Status} \\
        \hline
        \multirow{1}{*}{\textbf{Liquidity}} & Current Ratio & 1.72 & 1.50 & \textcolor{green}{Above} \\
        \hline
        \multirow{3}{*}{\textbf{Profitability}} & Gross Profit Ratio & 35.2\% & 32.0\% & \textcolor{green}{Above} \\
        & Operating Profit Ratio & 11.8\% & 10.5\% & \textcolor{green}{Above} \\
        & Net Profit Ratio & 7.25\% & 6.5\% & \textcolor{green}{Above} \\
        \hline
        \multirow{3}{*}{\textbf{Returns}} & ROE & 17.12\% & 15.0\% & \textcolor{green}{Above} \\
        & ROCE & 22.45\% & 18.0\% & \textcolor{green}{Above} \\
        & ROA & 12.02\% & 10.0\% & \textcolor{green}{Above} \\
        \hline
        \textbf{Solvency} & Debt Equity Ratio & 0.08 & 0.40 & \textcolor{green}{Better} \\
        \hline
        \multirow{3}{*}{\textbf{Valuation}} & EPS & \rupee~20.89 & -- & Growing \\
        & DPS & \rupee~5.50 & -- & Growing \\
        & P/E Ratio & 62.5x & 45.0x & Premium \\
        \hline
    \end{tabular}
\end{table}

%----------------------------------------------------------------------
\subsection{Financial Performance Trends}
%----------------------------------------------------------------------

\begin{figure}[H]
    \centering
    \begin{tabular}{|l|c|c|c|c|c|}
        \hline
        \textbf{Parameter} & \textbf{FY 18-19} & \textbf{FY 19-20} & \textbf{FY 20-21} & \textbf{FY 21-22} & \textbf{FY 22-23} \\
        \hline
        Revenue (Cr) & 9,636 & 9,587 & 10,328 & 14,856 & 18,063 \\
        \hline
        PAT (Cr) & 726 & 664 & 942 & 1,184 & 1,310 \\
        \hline
        EPS (\rupee) & 11.58 & 10.59 & 15.02 & 18.89 & 20.89 \\
        \hline
        ROE (\%) & 16.8 & 14.5 & 19.9 & 17.3 & 17.1 \\
        \hline
    \end{tabular}
    \caption{5-Year Financial Trends}
\end{figure}

\subsection{Key Financial Observations}

\begin{enumerate}[leftmargin=*]
    \item \textbf{Revenue Growth:} 87\% revenue growth over 5 years (CAGR ~13.4\%)
    \item \textbf{Profitability:} Consistent margins despite raw material inflation
    \item \textbf{Capital Structure:} Almost debt-free, ensuring financial flexibility
    \item \textbf{Shareholder Returns:} Consistent dividend payouts with growing EPS
    \item \textbf{Market Valuation:} Premium valuation reflecting growth potential and brand strength
    \item \textbf{Asset Efficiency:} Superior return ratios compared to industry peers
\end{enumerate}

%======================================================================
% CONCLUSION
%======================================================================
\newpage
\section{Conclusion}

Havells India Ltd. demonstrates excellence across all five functions of management:

\begin{enumerate}[leftmargin=*]
    \item \textbf{Planning:} Clear vision, well-defined objectives, and robust strategic planning for sustainable growth
    
    \item \textbf{Organising:} Well-structured organization with clear departmentalisation and authority flow
    
    \item \textbf{Staffing:} Comprehensive HR practices including recruitment, training, and competitive compensation
    
    \item \textbf{Directing:} Effective leadership, motivation techniques, and multi-channel communication
    
    \item \textbf{Controlling:} Rigorous performance standards and timely corrective actions
\end{enumerate}

Operating in an oligopolistic market, Havells has established itself as a premium brand with strong competitive advantages. The company's financial performance is exemplary, with:
\begin{itemize}
    \item Strong liquidity and minimal leverage
    \item Profitability ratios above industry averages
    \item Superior return on equity and capital employed
    \item Growing EPS and consistent dividends
    \item Premium market valuation reflecting investor confidence
\end{itemize}

Havells India Ltd. stands as a model of effective management practices in the Indian consumer goods industry.

%======================================================================
% REFERENCES
%======================================================================
\newpage
\section{References}

\begin{enumerate}[leftmargin=*]
    \item Havells India Ltd. Annual Report 2022-23
    \item Havells India Ltd. Annual Report 2021-22
    \item Havells India Ltd. Official Website: \url{https://www.havells.com}
    \item BSE/NSE Corporate Filings - Havells India Ltd.
    \item Company Investor Presentations (Q4 FY23)
    \item CRISIL Industry Reports - Electrical Equipment Industry
    \item Ministry of Corporate Affairs Filings
\end{enumerate}

%======================================================================
% END OF DOCUMENT
%======================================================================
\end{document}
